\section{SCENARIO \#9: 'Second Kernstown'': July 23-28, 1864
}


\begin{scenariodata}
\textbf{Game Length:} 6 Turns\\
\end{scenariodata}

\begin{scenariodata}
(Weather: July - October Column)\\
\end{scenariodata}

\begin{scenariodata}
\textbf{First Phase:} Confederates\\
\end{scenariodata}

\begin{scenariodata}
\textbf{Union:}\\
\end{scenariodata}

\begin{scenariodata}
Resource Factor 26\\
\end{scenariodata}

\begin{scenariodata}
\textbf{Initial Placement:}\\
\end{scenariodata}

\begin{scenariodata}
Hex N-22 (Kernstown) : 241, 3A, 6C, 1H,\\
\end{scenariodata}

\begin{scenariodata}
\textbf{and adjacent :} (Crook, 2F, 3S, 3W)\\
\end{scenariodata}

\begin{scenariodata}
B\&O : [151, 3A, 4C], (Hunter, 3F, 2S)\\
\end{scenariodata}

\begin{scenariodata}
\textbf{Withdrawal Hexes:}\\
\end{scenariodata}

\begin{scenariodata}
N-7, 8-10, 2-13, CC-15 through CC-33\\
(Supply/Wagon Entry Hexes:\\
\end{scenariodata}

\begin{scenariodata}
N-7, S-10, Z-13, A-11, CC-15 through CC-33\\
\end{scenariodata}

Three additional Fortification counters available.

\begin{scenariodata}
\textbf{Confederates:}\\
\end{scenariodata}

\begin{scenariodata}
\textbf{Resource Factor :} 10\\
\end{scenariodata}

\begin{scenariodata}
\textbf{Initial Placement:}\\
\end{scenariodata}

\begin{scenariodata}
Hex M-26 (Strasburg) : 401, 10A, 7C, 2H,\\
\end{scenariodata}

\begin{scenariodata}
\textbf{and adjacent :} (Early, Rodes, Gordon,\\
Breckinridge, Ramseur,\\
5S, 6W)\\
\end{scenariodata}

\begin{scenariodata}
Hex 0-51 (Staunton) : 11, (1F)\\
\end{scenariodata}

\begin{scenariodata}
Hex BB-53 (Charlottesville) +: 21, 1A, (1F, 1S)\\
\textbf{Partisans:} Mosby, Gilmor, McNeill\}\\
\end{scenariodata}

\begin{scenariodata}
\textbf{Withdrawal Hexes:}\\
\end{scenariodata}

\begin{scenariodata}
CC-40 through CC-58\\
\end{scenariodata}

\begin{scenariodata}
(Supply/Wagon Entry Hexes:\\
\end{scenariodata}

\begin{scenariodata}
CC-40 through CC-58, O-51, M-50, E-46)\\
\end{scenariodata}

Three additional Fortification counters available.

\begin{scenariodata}
\textbf{SPECIAL:} The Confederates get two Victory Points for each regular\\
\end{scenariodata}

Cavalry and/or Horse Artillery Combat Factor that exits the north edge of the Mapboard from Hex rows K through T, inclusive.

\begin{scenariodata}
\textbf{Victory Conditions:} Al! other results are a draw.\\
\end{scenariodata}

\begin{scenariodata}
\textbf{ADVANCED GAME:} Union wins if they have more total Victory Points\\
\end{scenariodata}

than the Confederates.

\begin{scenariodata}
Confederates win if they have 20 or more lead in\\
\end{scenariodata}

total Victory Points.

\begin{scenariodata}
\textbf{BASIC GAME:} Same as the Advanced Game, only "spot" the Con-\\
\end{scenariodata}

federates 5 Victory Points at the start of the game.

\begin{scenariodata}
\textbf{HISTORICAL RESULTS:} Confederate victory.\\
\end{scenariodata}

PLAYTESTER''S COMMENTS: This is a tough one for the Union to win. The Union Player should withdraw northwards as soon as possible (to the vicinity of Bunker Hill) and dig in, supported by part of the B\&O garrison in the Marginsburg area. The majority of the B\&O garrison should be stationed at Harper's Ferry. The Union must try to get some favorable attrition, withdrawing gradually, and staying between the Confederate cavalry and the edge of the Mapboard. This is one scenario where the Confederates can almost fight as much as they like. Drive forward as fast and as hard as possible, heading generally for the vicinity of Harper's Ferry.

The Confederates badly defeated the weak Union forces around Kernstown. Again a Rebel! army was able to advance to the Potomac and cut the B\&O Railroad. Confederate cavalry was sent north of the Potomac for a raid into Maryland. The town of Chambersburg was burned by the Confederate raiders, an action that brought cries of outrage from throughout the North. U.S. Grant reacted by deciding to settle the issue in the Shenandoah Valley once and for all. He dispatched further troops to the area, and placed all units in the region, for the first time, in one unified command, under his most fiery subordinate, Sheridan.

At this point, it will be useful to examine the two armies that would contest the control of the Shenandoah for three bloody months. The Confederate army was led by Jubal A. Early (1816-1894), a profane tobacco-chewing bachelor with an excellent combat record, who was regarded throughout the South at this time as a second "Stonewall" Jackson. Early's chief faults were his over-confidence, which made him too eager to fight the superior Union forces, and his faulty use of his cavalry, which never approached the effectiveness it had known under Turner Ashby. Early's subordinates were an extremely able, and distinguished group. Breckinridge commanded a temporary corps that consisted largely of his own division (the victors of New Market), and Gordon's division of the Il Corps. Rodes commanded another temporary corps consisting of his own and Ramseur's divisions.

Robert E, Rodes, a graduate of VMI, was a mustachioed veteran of 35 who had spearheaded Jackson's last attack at Chancellorsville in 1863. John B. Gordon (1832-1904), a lawyer by profession, was an extremely able citizen-soldier and a savage fighter, who would rise to the command of a corps by the end of the war. Stephen D. Ramseur was, at 27, the youngest Major General in the Confederate army. He was a graduate of West Point (1860), and already carried the scars of three wounds, The division of Joseph B. Kershaw, a calm, deliberate veteran of 42, was an experienced outfit that was shuffled in and out of the Valley as the campaign progressed. The Confederate tl Corps, Army of Northern Virginia, that formed the core of Early's army, was, of course, "Stonewall" Jackson's old command and still contained many veterans of the 1862 Valley Campaign. These hardened veterans proved able to bounce back after repeated defeats, and had truly amazing morale (A II Corps officer noted a man in his ranks with both arms so wounded that he could not carry a rifle. Asked what he was doing on the firing line if he couldn't fight, the soldier replied, "1 know | can't fight, but I can still yell.''\},

The Union forces, now styled the Army of the Shenandoah, were Jed by Phillip H. Sheridan (1831-1888), a West Pointer (1853) who had built up a reputation as a fiery leader in command of an infantry division in the Army of the Cumberland, and, more recently, in command of the Cavalry Corps of the Army of the Potomac. A natural leader of men, Sheridan was also an excellent tactician. Horatio G. Wright (1820-1899) commanded the VI Corps, a veteran outfit from the Army of the Potomac that provided the best troops in the Union army. Wright was a fine corps commander, although his earlier experiences had shown that he was not fitted for an independent command. George Crook (1829-1890), later to become a renowned Indian fighter, was the commander of the VIIL Corps (also known as the Army of West Virginia), The VIII Corps was made up largely of the defeated remnants of earlier Union Valley armies, and had an uncertain morale, although, in general, they performed well under Crook. William H. Emory, a much older officer than the others (West Point, 1831), led the XIX Corps, a unit that had been transferred from Louisiana. The XtX Corps performed well under Emory's steady leadership, but the troops were never really considered to be the equal of the VI Corps. The Union cavalry was excellent, being largely composed of units from Sheridan's Cavalry Corps.

Sheridan's army had been rather hastily thrown together, and was not yet completely organized. Early and his forces had been together for some time, and were well-practiced in battle and maneuver. The clamor for action from both press and public finally forced the Union forces to move forward to test the Confederate strength...

